\documentclass[11pt]{amsart}

\usepackage[a4paper,margin=1in]{geometry}
\usepackage[T1]{fontenc}
\usepackage{lmodern}
\usepackage{microtype}
\usepackage{amsmath,amssymb,amsthm,mathtools}
\usepackage[colorlinks=true,linkcolor=blue,citecolor=blue,urlcolor=blue]{hyperref}

\numberwithin{equation}{section}

\DeclareMathOperator{\tridiag}{tridiag}
\DeclareMathOperator{\diag}{diag}
\DeclareMathOperator{\spec}{spec}
\DeclareMathOperator{\dist}{dist}

\newcommand{\C}{\mathbb C}
\newcommand{\R}{\mathbb R}
\newcommand{\eps}{\varepsilon}
\newcommand{\ii}{\mathrm i}
\newcommand{\A}{A}
\newcommand{\B}{B}
\newcommand{\I}{I}
\newcommand{\J}{J}
\newcommand{\HH}{H}

\newtheorem{theorem}{Theorem}[section]
\newtheorem{mainthm}[theorem]{Main Theorem}
\newtheorem{lemma}[theorem]{Lemma}
\newtheorem{proposition}[theorem]{Proposition}
\newtheorem{corollary}[theorem]{Corollary}
\newtheorem{definition}[theorem]{Definition}
\newtheorem{conjecture}[theorem]{Conjecture}
\newtheorem{remark}[theorem]{Remark}

\title[Connectedness thresholds for Toeplitz pseudospectra]
{Real-Axis Barriers and Connectedness Thresholds for Pseudospectra of a Nonnormal Tridiagonal Toeplitz Family}

\author{Anonymous}

\subjclass[2020]{15A18, 15B05, 47A10, 65F15, 81Q12}
\keywords{Pseudospectrum, nonnormal matrix, tridiagonal Toeplitz matrix, least singular value, connectedness, Hatano--Nelson model, non-Hermitian skin effect}

\begin{document}

\begin{abstract}
We study the Euclidean pseudospectra of the finite nonnormal tridiagonal Toeplitz matrices
\[
A_n(a)=\tridiag(a,0,1),\qquad 0<a<1,\qquad n\ge2.
\]
The spectrum is real because \(A_n(a)\) is diagonally similar to a symmetric Jacobi matrix, but its Euclidean pseudospectra are genuinely nonnormal.  We introduce the real-axis barrier
\[
\tau_n(a)
=
\max_{1\le k\le n-1}
\max_{\lambda_{n,k}(a)\le x\le \lambda_{n,k+1}(a)}
s_{\min}(xI_n-A_n(a)),
\]
where \(\lambda_{n,1}(a)<\cdots<\lambda_{n,n}(a)\) are the eigenvalues of \(A_n(a)\).  The barrier gives the exact connectedness criterion
\[
\sigma_\eps(A_n(a))\ \text{is connected}
\quad\Longleftrightarrow\quad
\eps>\tau_n(a).
\]
The key point is the vertical monotonicity theorem
\[
s_{\min}((x+\ii y)I_n-A_n(a))
\ge
s_{\min}(xI_n-A_n(a)),
\qquad x,y\in\R,
\]
which rules out off-real bridges before the corresponding real gap closes.  We also give explicit two-sided bounds for \(\tau_n(a)\).  The lower bound is expressed through the optimal moment naturally produced by the Green-kernel estimate,
\[
S_2(a)=\sum_{m=1}^{\infty}m^2a^m=\frac{a(1+a)}{(1-a)^3}.
\]
Consequently,
\[
c_2(a)(n+1)a^{n/2}
\le
\tau_n(a)
\le
(n+1)a^{n/2},
\qquad
c_2(a)=\frac{\sqrt8(1-a)^3}{3\pi^2(1+a)}.
\]
These estimates imply sharp logarithmic bounds for the first-hit connectedness threshold.  If the finite real-axis barrier monotonicity conjecture \(\tau_{n+1}(a)\le\tau_n(a)\) holds, then the first-hit threshold is also the eventual connectedness threshold.  Finally, we explain the interpretation for the open-boundary Hatano--Nelson chain.
\end{abstract}

\maketitle

\section{Introduction}

For \(M\in\C^{n\times n}\), its \(\eps\)-pseudospectrum is
\[
\sigma_\eps(M)
=
\left\{
z\in\C:\|(zI_n-M)^{-1}\|_2>\eps^{-1}
\right\},
\]
with the convention that the resolvent norm is \(+\infty\) on \(\sigma(M)\).  Equivalently,
\[
\sigma_\eps(M)
=
\left\{
z\in\C:s_{\min}(zI_n-M)<\eps
\right\}.
\]
Pseudospectra are central in nonnormal matrix analysis because eigenvalues alone may fail to describe resolvent growth and spectral instability; see Trefethen \cite{Trefethen1997} and Trefethen--Embree \cite{TrefethenEmbree2005}.

This article concerns
\[
A_n(a)
=
\tridiag(a,0,1)
=
\begin{bmatrix}
0 & 1 \\
a & 0 & \ddots \\
& \ddots & \ddots & \ddots \\
&& \ddots & 0 & 1 \\
&&& a & 0
\end{bmatrix},
\qquad 0<a<1.
\]
Although \(A_n(a)\) is nonnormal, it is diagonally similar to the symmetric matrix
\[
\sqrt a\,\tridiag(1,0,1).
\]
Hence its eigenvalues are real and explicit.  However, pseudospectra in the Euclidean norm are not invariant under nonunitary similarity.  Thus the topology of \(\sigma_\eps(A_n(a))\) reflects nonnormality.

For fixed \(a\in(0,1)\), define the real-axis least singular value
\[
\mu_n(x;a):=s_{\min}(xI_n-A_n(a)),
\qquad x\in\R.
\]
Let
\[
\lambda_{n,1}(a)<\lambda_{n,2}(a)<\cdots<\lambda_{n,n}(a)
\]
be the eigenvalues of \(A_n(a)\).  We define the \(k\)-th gap height by
\[
\beta_{n,k}(a)
:=
\max_{\lambda_{n,k}(a)\le x\le\lambda_{n,k+1}(a)}
\mu_n(x;a),
\qquad k=1,\dots,n-1,
\]
and the largest real-axis barrier by
\[
\tau_n(a):=\max_{1\le k\le n-1}\beta_{n,k}(a).
\]
This quantity is the obstruction to connectedness along the real axis.  We prove below that vertical monotonicity reduces planar connectedness of \(\sigma_\eps(A_n(a))\) exactly to connectedness of its real trace.

The desired finite-dimensional monotonicity statement is
\[
\tau_{n+1}(a)\le\tau_n(a),
\qquad 0<a<1.
\]
If true, it immediately implies that connectedness, once it appears, never disappears as \(n\) increases.  Because this comparison is the remaining hard finite-\(n\) point, we separate the unconditional connectedness criterion and estimates from the conditional eventual-threshold consequence.

\section{Elementary structure}

Let
\[
B_n(z):=zI_n-A_n(a),
\qquad
s_n(z):=s_{\min}(B_n(z)).
\]
Let \(J_n\) be the reversal matrix,
\[
J_ne_j=e_{n+1-j}.
\]
Then
\[
A_n(a)^TJ_n=J_nA_n(a).
\]
Consequently, for real \(x\),
\[
H_n(x):=J_nB_n(x)=J_n(xI_n-A_n(a))
\]
is real symmetric.  Since \(J_n\) is orthogonal,
\begin{equation}\label{eq:real-axis-smin}
s_n(x)
=
s_{\min}(H_n(x))
=
\min_{\rho\in\sigma(H_n(x))}|\rho|.
\end{equation}

The spectrum of \(A_n(a)\) is explicit.  Put
\[
\Delta_n:=\diag(1,a^{1/2},a,\dots,a^{(n-1)/2}).
\]
Then
\[
\Delta_n^{-1}A_n(a)\Delta_n
=
\sqrt a\,\tridiag(1,0,1).
\]
Therefore
\begin{equation}\label{eq:toeplitz-spectrum}
\sigma(A_n(a))
=
\left\{
2\sqrt a\cos\frac{j\pi}{n+1}:j=1,\dots,n
\right\}.
\end{equation}
We use the increasing order
\[
\lambda_{n,k}(a)
=
2\sqrt a\cos\frac{(n+1-k)\pi}{n+1},
\qquad k=1,\dots,n.
\]
Thus
\[
\lambda_{n,1}(a)<\lambda_{n,2}(a)<\cdots<\lambda_{n,n}(a).
\]

\section{The real-axis barrier}

\begin{definition}
For \(n\ge2\), \(0<a<1\), and \(1\le k\le n-1\), set
\[
G_{n,k}(a)
:=
[\lambda_{n,k}(a),\lambda_{n,k+1}(a)].
\]
Define
\[
\beta_{n,k}(a)
:=
\max_{x\in G_{n,k}(a)}
s_{\min}(xI_n-A_n(a)),
\]
and
\[
\tau_n(a)
:=
\max_{1\le k\le n-1}\beta_{n,k}(a).
\]
\end{definition}

The maximum exists because \(x\mapsto s_{\min}(xI_n-A_n(a))\) is continuous and each \(G_{n,k}(a)\) is compact.  Since \(xI_n-A_n(a)\) is singular exactly at the endpoints of \(G_{n,k}(a)\), one has
\[
\beta_{n,k}(a)>0.
\]

\begin{proposition}[Real-axis connectedness criterion]\label{prop:real-axis-connectedness}
Let
\[
E_{n,\eps}(a)
:=
\{x\in\R:s_{\min}(xI_n-A_n(a))<\eps\}.
\]
Then
\[
E_{n,\eps}(a)\ \text{is connected}
\quad\Longleftrightarrow\quad
\eps>\tau_n(a).
\]
\end{proposition}

\begin{proof}
The function
\[
x\mapsto\mu_n(x;a):=s_{\min}(xI_n-A_n(a))
\]
is continuous and nonnegative.  Its zeros are precisely the real eigenvalues
\[
\lambda_{n,1}(a),\dots,\lambda_{n,n}(a).
\]
Therefore \(E_{n,\eps}(a)\) contains a neighborhood of every eigenvalue.

The only possible disconnections of \(E_{n,\eps}(a)\) between eigenvalues occur inside the closed gaps
\[
G_{n,k}(a)=[\lambda_{n,k}(a),\lambda_{n,k+1}(a)].
\]
The \(k\)-th gap is fully filled by the strict sublevel set \(\mu_n<\eps\) if and only if
\[
\mu_n(x;a)<\eps
\qquad
\text{for every }x\in G_{n,k}(a).
\]
Since
\[
\beta_{n,k}(a)=\max_{x\in G_{n,k}(a)}\mu_n(x;a),
\]
this is equivalent to
\[
\eps>\beta_{n,k}(a).
\]
All real gaps are filled if and only if
\[
\eps>\max_{1\le k\le n-1}\beta_{n,k}(a)=\tau_n(a).
\]
The strict inequality is essential because the pseudospectrum is defined by \(s_{\min}<\eps\).  If \(\eps=\tau_n(a)\), then at least one real gap contains a point where \(\mu_n=\eps\), and that point is not included in \(E_{n,\eps}(a)\).
\end{proof}

\section{Vertical monotonicity and planar connectedness}
\label{sec:vertical-planar}

The preceding real-axis criterion becomes a planar connectedness criterion once one knows that the least singular value increases along vertical lines.  In the present Toeplitz family this property is true.

\begin{theorem}[Vertical monotonicity]\label{thm:vertical}
For every \(0<a<1\), every \(n\ge2\), and every \(x,y\in\R\),
\[
s_{\min}((x+\ii y)I_n-A_n(a))
\ge
s_{\min}(xI_n-A_n(a)).
\]
Equivalently, for every fixed \(x\in\R\), the map
\[
[0,\infty)\ni y
\longmapsto
s_{\min}((x+\ii y)I_n-A_n(a))
\]
is nondecreasing.
\end{theorem}

The proof is given in a sequence of lemmas.  Throughout the proof we fix \(0<a<1\) and \(n\ge2\), and write
\[
A:=A_n(a),
\qquad
B_{x,y}:=(x+\ii y)I_n-A,
\qquad
\mathcal H_x(y):=B_{x,y}^*B_{x,y}.
\]
We also write
\[
s(x,y):=s_{\min}(B_{x,y}),
\qquad
\lambda_1(x,y):=s(x,y)^2
=
\lambda_{\min}(\mathcal H_x(y)).
\]
Let \(J=J_n\) be the reversal matrix.  Then
\[
JAJ=A^{\mathsf T}.
\]

\begin{lemma}[The case \(n=2\)]\label{lem:vertical-n2}
For \(n=2\) and all \(x\in\R\), the function
\[
y\longmapsto \lambda_1(x,y)
\]
is nondecreasing on \([0,\infty)\).
\end{lemma}

\begin{proof}
For \(n=2\),
\[
B_{x,y}
=
\begin{bmatrix}
x+\ii y & -1\\
-a & x+\ii y
\end{bmatrix},
\]
and hence
\[
\mathcal H_x(y)
=
\begin{bmatrix}
a^2+x^2+y^2 & -(1+a)x+\ii(1-a)y\\
-(1+a)x-\ii(1-a)y & 1+x^2+y^2
\end{bmatrix}.
\]
Therefore
\[
\lambda_1(x,y)
=
x^2+y^2+\frac{1+a^2}{2}
-
\frac12
\sqrt{(1-a^2)^2+4(1+a)^2x^2+4(1-a)^2y^2}.
\]
Differentiation gives
\[
\partial_y\lambda_1(x,y)
=
2y
\left(
1-
\frac{(1-a)^2}
{\sqrt{(1-a^2)^2+4(1+a)^2x^2+4(1-a)^2y^2}}
\right).
\]
Since
\[
\sqrt{(1-a^2)^2+4(1+a)^2x^2+4(1-a)^2y^2}
\ge
1-a^2
\ge
(1-a)^2,
\]
we have
\[
\partial_y\lambda_1(x,y)\ge0
\qquad(y\ge0).
\]
\end{proof}

For the rest of the proof of Theorem \ref{thm:vertical}, assume \(n\ge3\) and \(y>0\).  Put
\[
z:=x+\ii y,
\qquad
\lambda:=\lambda_1(x,y),
\qquad
c:=\overline z+az=(1+a)x-\ii(1-a)y.
\]
Define
\[
d_-:=a^2+|z|^2-\lambda,
\qquad
d:=1+a^2+|z|^2-\lambda,
\qquad
d_+:=1+|z|^2-\lambda.
\]
Then
\[
M:=\mathcal H_x(y)-\lambda I_n
\]
is the Hermitian pentadiagonal matrix with first diagonal entry \(d_-\), last diagonal entry \(d_+\), interior diagonal entries \(d\), first off-diagonal entries \(-c\) above the diagonal and \(-\overline c\) below the diagonal, and second off-diagonal entries \(a\).

For an index set \(I\subset\{1,\dots,n\}\), write \(M[I]\) for the principal submatrix of \(M\) indexed by \(I\).  For \(2\le j\le n-1\), set
\[
L_j:=M[\{1,\dots,j-1\}],
\qquad
R_j:=M[\{j+1,\dots,n\}].
\]

\begin{lemma}[One-sided positivity]\label{lem:one-sided-positivity}
For every \(2\le j\le n-1\),
\[
L_j\succ0,
\qquad
R_j\succ0.
\]
In particular,
\[
d_->0,
\qquad
d_+>0.
\]
\end{lemma}

\begin{proof}
Since \(\lambda=\lambda_{\min}(\mathcal H_x(y))\), one has \(M\succeq0\).  Hence every principal submatrix of \(M\) is positive semidefinite.  It remains to show that the kernels of \(L_j\) and \(R_j\) are trivial.

We prove this for \(R_j\); the proof for \(L_j\) is the same after reversing the indices.  Suppose \(w\in\ker R_j\), and extend \(w\) to a vector \(\widehat w\in\C^n\) by placing zeros in the first \(j\) positions.  The rows \(j+1,\dots,n\) of \(M\widehat w\) vanish.  Row \(j-1\) gives
\[
aw_1=0,
\]
so \(w_1=0\).  If \(w\) has a second component, row \(j\) then gives
\[
aw_2=0,
\]
so \(w_2=0\).  The first row of \(R_jw=0\) now gives \(aw_3=0\), if such a component exists.  Continuing in this way gives \(w=0\).  Thus \(R_j\succ0\).  Similarly, \(L_j\succ0\).

Finally,
\[
L_2=[d_-],
\qquad
R_{n-1}=[d_+],
\]
so \(d_->0\) and \(d_+>0\).
\end{proof}

Define
\[
T:=
\begin{bmatrix}
d & -c\\
-\overline c & d
\end{bmatrix},
\qquad
E:=
\begin{bmatrix}
a & 0\\
-c & a
\end{bmatrix}.
\]
For \(3\le j\le n\), let \(S_j^{\rm L}\) be the Schur complement of \(M[\{1,\dots,j-1\}]\) onto its last two coordinates.  For \(1\le j\le n-2\), let \(S_j^{\rm R}\) be the Schur complement of \(M[\{j+1,\dots,n\}]\) onto its first two coordinates.  If the eliminated block is empty, the Schur complement is understood as the corresponding \(2\times2\) matrix itself.

\begin{lemma}[Schur-complement half-plane property]\label{lem:schur-half-plane}
Whenever
\[
S_j^{\rm L}
=
\begin{bmatrix}
\alpha_j^{\rm L} & \beta_j^{\rm L}\\
\overline{\beta_j^{\rm L}} & \delta_j^{\rm L}
\end{bmatrix},
\]
one has
\[
\Im\beta_j^{\rm L}>0.
\]
Whenever
\[
S_j^{\rm R}
=
\begin{bmatrix}
\alpha_j^{\rm R} & \beta_j^{\rm R}\\
\overline{\beta_j^{\rm R}} & \delta_j^{\rm R}
\end{bmatrix},
\]
one has
\[
\Im\beta_j^{\rm R}>0.
\]
\end{lemma}

\begin{proof}
We first record the elementary propagation rule.  Let
\[
S=
\begin{bmatrix}
\alpha & \beta\\
\overline\beta & \delta
\end{bmatrix}
\succ0,
\qquad
\Delta:=\alpha\delta-|\beta|^2>0.
\]
A direct calculation gives
\[
T-E^*S^{-1}E
=
\begin{bmatrix}
\alpha' & \beta'\\
\overline{\beta'} & \delta'
\end{bmatrix},
\]
where
\[
\beta'
=
\frac{a^2\beta+a\alpha\,\overline c-c\Delta}{\Delta}.
\]
Since
\[
\Im c=-(1-a)y,
\qquad
\Im\overline c=(1-a)y,
\]
we obtain
\[
\Im\beta'
=
\frac{
a^2\Im\beta+a\alpha(1-a)y+\Delta(1-a)y
}{\Delta}.
\]
Thus \(\Im\beta>0\) implies \(\Im\beta'>0\).

Similarly,
\[
T-ES^{-1}E^*
=
\begin{bmatrix}
\widetilde\alpha & \widetilde\beta\\
\overline{\widetilde\beta} & \widetilde\delta
\end{bmatrix},
\]
where
\[
\widetilde\beta
=
\frac{a^2\beta+a\delta\,\overline c-c\Delta}{\Delta},
\]
and therefore
\[
\Im\widetilde\beta
=
\frac{
a^2\Im\beta+a\delta(1-a)y+\Delta(1-a)y
}{\Delta}.
\]
Again \(\Im\beta>0\) implies \(\Im\widetilde\beta>0\).

For the left Schur complements, the initial cases are
\[
S_3^{\rm L}
=
\begin{bmatrix}
d_- & -c\\
-\overline c & d
\end{bmatrix},
\]
and, when meaningful,
\[
S_4^{\rm L}
=
\begin{bmatrix}
d-\dfrac{|c|^2}{d_-}
&
-c+\dfrac{a\overline c}{d_-}
\\[2mm]
-\overline c+\dfrac{ac}{d_-}
&
d-\dfrac{a^2}{d_-}
\end{bmatrix}.
\]
Both have upper-right entry with positive imaginary part, because \(d_->0\) and \(y>0\).  For \(j\ge5\), block Schur-complement calculation gives
\[
S_j^{\rm L}
=
T-E^*(S_{j-2}^{\rm L})^{-1}E.
\]
The propagation rule proves \(\Im\beta_j^{\rm L}>0\) for all \(3\le j\le n\).

For the right Schur complements, the terminal initial cases are
\[
S_{n-2}^{\rm R}
=
\begin{bmatrix}
d & -c\\
-\overline c & d_+
\end{bmatrix},
\]
and, when meaningful,
\[
S_{n-3}^{\rm R}
=
\begin{bmatrix}
d-\dfrac{a^2}{d_+}
&
-c+\dfrac{a\overline c}{d_+}
\\[2mm]
-\overline c+\dfrac{ac}{d_+}
&
d-\dfrac{|c|^2}{d_+}
\end{bmatrix}.
\]
Again the upper-right entries have positive imaginary part.  For \(j\le n-4\),
\[
S_j^{\rm R}
=
T-E(S_{j+2}^{\rm R})^{-1}E^*.
\]
The second propagation rule proves \(\Im\beta_j^{\rm R}>0\) for all \(1\le j\le n-2\).
\end{proof}

\begin{lemma}[No interior zero]\label{lem:no-interior-zero}
Let
\[
v=(v_1,\dots,v_n)^T\in\ker M\setminus\{0\}.
\]
Then
\[
v_j\ne0,
\qquad
j=2,\dots,n-1.
\]
\end{lemma}

\begin{proof}
We first show that
\[
v_2\ne0,
\qquad
v_{n-1}\ne0.
\]
By reversal symmetry it is enough to prove \(v_2\ne0\).  Suppose \(v_2=0\).

If \(n=3\), then \(v_1\ne0\), because otherwise the positive definiteness of \(R_2=[d_+]\) would force \(v_3=0\).  Row \(3\) of \(Mv=0\) gives
\[
d_+v_3+av_1=0,
\]
so
\[
v_3=-\frac{a}{d_+}v_1.
\]
Row \(2\) then gives
\[
-\overline c\,v_1-cv_3=0,
\]
hence
\[
\overline c=\frac{a}{d_+}c.
\]
This is impossible because \(a/d_+>0\), whereas
\[
\Im\overline c=(1-a)y>0,
\qquad
\Im c=-(1-a)y<0.
\]

Now assume \(n\ge4\).  If \(v_1=0\), then
\[
R_2(v_3,\dots,v_n)^T=0,
\]
contradicting \(R_2\succ0\).  Hence \(v_1\ne0\).  Let
\[
G_2:=R_2^{-1},
\qquad
g:=\langle e_1,G_2e_1\rangle>0,
\qquad
h:=\langle e_2,G_2e_1\rangle,
\]
where \(e_1,e_2\) are the first two standard basis vectors in \(\C^{n-2}\).  From the equations on the coordinates \(3,\dots,n\), we get
\[
(v_3,\dots,v_n)^T=-av_1G_2e_1.
\]
Thus
\[
v_3=-agv_1,
\qquad
v_4=-ahv_1.
\]
Row \(1\) gives
\[
d_-v_1+av_3=0,
\]
so
\[
d_-=a^2g.
\]
Row \(2\) gives
\[
-\overline c\,v_1-cv_3+av_4=0,
\]
and therefore
\[
-\overline c+agc-a^2h=0.
\]
Equivalently,
\[
-c+a\frac{h}{g}
=
-\frac{1}{ag}\,\overline c.
\]
The right-hand side has negative imaginary part.

On the other hand, write
\[
S_2^{\rm R}
=
\begin{bmatrix}
\alpha_2^{\rm R} & \beta_2^{\rm R}\\
\overline{\beta_2^{\rm R}} & \delta_2^{\rm R}
\end{bmatrix}.
\]
The upper-left \(2\times2\) block of \(R_2^{-1}\) is \((S_2^{\rm R})^{-1}\), and hence
\[
\frac{h}{g}
=
-\frac{\overline{\beta_2^{\rm R}}}{\delta_2^{\rm R}}.
\]
Therefore
\[
-c+a\frac{h}{g}
=
-c-a\frac{\overline{\beta_2^{\rm R}}}{\delta_2^{\rm R}}.
\]
By Lemma \ref{lem:schur-half-plane},
\[
\Im\beta_2^{\rm R}>0,
\]
so the last expression has positive imaginary part.  This contradiction proves \(v_2\ne0\), and by reversal \(v_{n-1}\ne0\).

It remains to rule out zeros at indices \(3,\dots,n-2\).  Suppose that
\[
v_j=0
\]
for some \(3\le j\le n-2\).  Set
\[
G_j^{\rm L}:=L_j^{-1},
\qquad
G_j^{\rm R}:=R_j^{-1}.
\]
Let
\[
g_j^{\rm L}
:=
\langle e_{j-1},G_j^{\rm L}e_{j-1}\rangle>0,
\qquad
h_j^{\rm L}
:=
\langle e_{j-2},G_j^{\rm L}e_{j-1}\rangle,
\]
where \(e_{j-2},e_{j-1}\) denote the last two basis vectors in \(\C^{j-1}\), and let
\[
g_j^{\rm R}
:=
\langle e_1,G_j^{\rm R}e_1\rangle>0,
\qquad
h_j^{\rm R}
:=
\langle e_2,G_j^{\rm R}e_1\rangle,
\]
where \(e_1,e_2\) denote the first two basis vectors in \(\C^{n-j}\).

Splitting \(v\) at the zero coordinate \(j\), the left and right equations give
\[
v_{j-1}=-ag_j^{\rm L}v_{j+1},
\qquad
v_{j+1}=-ag_j^{\rm R}v_{j-1}.
\]
Thus
\[
v_{j-1}\ne0,
\qquad
v_{j+1}\ne0,
\qquad
a^2g_j^{\rm L}g_j^{\rm R}=1.
\]
Moreover,
\[
v_{j-2}
=
\frac{h_j^{\rm L}}{g_j^{\rm L}}v_{j-1},
\qquad
v_{j+2}
=
\frac{h_j^{\rm R}}{g_j^{\rm R}}v_{j+1}.
\]
Row \(j\) of \(Mv=0\) is
\[
av_{j-2}-\overline c\,v_{j-1}-cv_{j+1}+av_{j+2}=0.
\]
Hence
\[
m_j^{\rm L}v_{j-1}+m_j^{\rm R}v_{j+1}=0,
\]
where
\[
m_j^{\rm L}
:=
a\frac{h_j^{\rm L}}{g_j^{\rm L}}-\overline c,
\qquad
m_j^{\rm R}
:=
-c+a\frac{h_j^{\rm R}}{g_j^{\rm R}}.
\]
Using \(v_{j-1}=-ag_j^{\rm L}v_{j+1}\), we obtain
\[
m_j^{\rm R}=ag_j^{\rm L}m_j^{\rm L}.
\]
Thus \(m_j^{\rm R}\) is a positive real multiple of \(m_j^{\rm L}\).

We now compute their imaginary parts.  Write
\[
S_j^{\rm L}
=
\begin{bmatrix}
\alpha_j^{\rm L} & \beta_j^{\rm L}\\
\overline{\beta_j^{\rm L}} & \delta_j^{\rm L}
\end{bmatrix},
\qquad
S_j^{\rm R}
=
\begin{bmatrix}
\alpha_j^{\rm R} & \beta_j^{\rm R}\\
\overline{\beta_j^{\rm R}} & \delta_j^{\rm R}
\end{bmatrix}.
\]
The corresponding inverse-block formulas give
\[
\frac{h_j^{\rm L}}{g_j^{\rm L}}
=
-\frac{\beta_j^{\rm L}}{\alpha_j^{\rm L}},
\qquad
\frac{h_j^{\rm R}}{g_j^{\rm R}}
=
-\frac{\overline{\beta_j^{\rm R}}}{\delta_j^{\rm R}}.
\]
Therefore
\[
m_j^{\rm L}
=
-\overline c-a\frac{\beta_j^{\rm L}}{\alpha_j^{\rm L}},
\qquad
m_j^{\rm R}
=
-c-a\frac{\overline{\beta_j^{\rm R}}}{\delta_j^{\rm R}}.
\]
By Lemma \ref{lem:schur-half-plane},
\[
\Im m_j^{\rm L}<0<\Im m_j^{\rm R}.
\]
This is incompatible with
\[
m_j^{\rm R}=ag_j^{\rm L}m_j^{\rm L},
\qquad
ag_j^{\rm L}>0.
\]
Hence no such \(j\) exists.
\end{proof}

\begin{lemma}[No boundary zero]\label{lem:no-boundary-zero}
Let
\[
v=(v_1,\dots,v_n)^T\in\ker M\setminus\{0\}.
\]
Then
\[
v_1\ne0,
\qquad
v_n\ne0.
\]
\end{lemma}

\begin{proof}
We prove \(v_1\ne0\); the proof of \(v_n\ne0\) is symmetric.

Assume \(v_1=0\).  By Lemma \ref{lem:no-interior-zero}, \(v_2\ne0\).  Let
\[
R_1:=M[\{2,\dots,n\}].
\]
Then
\[
(v_2,\dots,v_n)^T\in\ker R_1.
\]
Taking the Schur complement of \(R_1\) onto its first two coordinates gives
\[
S_1^{\rm R}
=
\begin{bmatrix}
\alpha & \beta\\
\overline\beta & \delta
\end{bmatrix},
\qquad
\Im\beta>0
\]
by Lemma \ref{lem:schur-half-plane}.  The vector \((v_2,v_3)^T\) belongs to \(\ker S_1^{\rm R}\).  Since \(S_1^{\rm R}\succeq0\) and \(\beta\ne0\), one has \(\alpha>0\).  From the first row,
\[
\alpha v_2+\beta v_3=0,
\]
so
\[
\frac{v_3}{v_2}=-\frac{\alpha}{\beta}.
\]
Because \(\alpha>0\) and \(\Im\beta>0\), this ratio has positive imaginary part.

On the other hand, row \(1\) of \(Mv=0\) gives
\[
-cv_2+av_3=0,
\]
so
\[
\frac{v_3}{v_2}=\frac{c}{a}.
\]
But
\[
\Im\frac{c}{a}
=
-\frac{(1-a)y}{a}<0,
\]
a contradiction.  Thus \(v_1\ne0\).

If \(v_n=0\), then \((v_1,\dots,v_{n-1})^T\in\ker M[\{1,\dots,n-1\}]\).  Taking the Schur complement onto the last two coordinates gives
\[
S_n^{\rm L}
=
\begin{bmatrix}
\alpha & \beta\\
\overline\beta & \delta
\end{bmatrix},
\qquad
\Im\beta>0.
\]
The vector \((v_{n-2},v_{n-1})^T\) belongs to \(\ker S_n^{\rm L}\).  Since \(v_{n-1}\ne0\) by Lemma \ref{lem:no-interior-zero}, the second row gives
\[
\frac{v_{n-2}}{v_{n-1}}
=
-\frac{\delta}{\overline\beta},
\]
which has negative imaginary part.  Row \(n\) of \(Mv=0\), however, gives
\[
av_{n-2}-\overline c\,v_{n-1}=0,
\]
so
\[
\frac{v_{n-2}}{v_{n-1}}
=
\frac{\overline c}{a},
\]
whose imaginary part is \((1-a)y/a>0\).  This contradiction proves \(v_n\ne0\).
\end{proof}

\begin{corollary}[Simplicity of the lowest branch]\label{cor:lowest-simple}
For \(n\ge3\), \(x\in\R\), and \(y>0\), the eigenvalue
\[
\lambda_1(x,y)=\lambda_{\min}(\mathcal H_x(y))
\]
is simple.
\end{corollary}

\begin{proof}
If the eigenspace had dimension at least \(2\), then there would be a nonzero ground-state vector \(v\in\ker M\) with \(v_2=0\).  This contradicts Lemma \ref{lem:no-interior-zero}.
\end{proof}

\begin{lemma}[Canonical singular-vector gauge]\label{lem:canonical-gauge}
Assume \(n\ge3\), \(y>0\), and that one variable is restricted to an interval on which \(\lambda_1(x,y)\) is simple.  Then one can choose real-analytic unit singular vectors \(u,v\in\C^n\) such that
\[
B_{x,y}v=s(x,y)u,
\qquad
B_{x,y}^*u=s(x,y)v,
\qquad
u=J\overline v.
\]
Consequently, with \(v_0=v_{n+1}=0\),
\begin{equation}\label{eq:canonical-coordinate}
(x+\ii y)v_j-av_{j-1}-v_{j+1}
=
s(x,y)\overline{v_{n+1-j}},
\qquad
j=1,\dots,n.
\end{equation}
\end{lemma}

\begin{proof}
Since the lowest eigenvalue of \(\mathcal H_x(y)\) is simple on the interval, analytic perturbation theory gives a real-analytic unit eigenvector \(v\) for \(\lambda_1(x,y)\).  Define
\[
u:=s(x,y)^{-1}B_{x,y}v.
\]
Then
\[
B_{x,y}v=s(x,y)u,
\qquad
B_{x,y}^*u=s(x,y)v.
\]

The identity \(JAJ=A^{\mathsf T}\) implies
\[
B_{x,y}J=JB_{x,y}^{\mathsf T},
\qquad
B_{x,y}^*J=J\overline{B_{x,y}}.
\]
Therefore, if \((u,v)\) is a singular pair, then
\[
(J\overline v,J\overline u)
\]
is also a singular pair for the same singular value, with the first vector left and the second vector right.  Since the singular value is simple, there is a unimodular scalar \(\alpha\) such that
\[
J\overline v=\alpha u.
\]
Multiplying both \(u\) and \(v\) by a unimodular square root of \(\alpha\), we obtain the gauge
\[
u=J\overline v.
\]
The coordinate equation \eqref{eq:canonical-coordinate} is just the \(j\)-th component of \(B_{x,y}v=s(x,y)J\overline v\).
\end{proof}

\begin{lemma}[Current identity]\label{lem:current-identity}
Assume \(n\ge3\), \(y>0\), and let \(v\) be in the canonical gauge of Lemma \ref{lem:canonical-gauge}.  Define
\[
\mathcal J_j:=\Im(\overline{v_j}v_{j+1}),
\qquad
j=1,\dots,n-1,
\]
and set
\[
\mathcal J_0=\mathcal J_n=0.
\]
Then
\begin{equation}\label{eq:current-identity}
\mathcal J_j-a\mathcal J_{j-1}
=
y|v_j|^2+s(x,y)\Im(v_jv_{n+1-j}),
\qquad
j=1,\dots,n,
\end{equation}
and
\begin{equation}\label{eq:lambda-y-current}
\partial_y\lambda_1(x,y)
=
2y-2(1-a)\sum_{j=1}^{n-1}\mathcal J_j.
\end{equation}
\end{lemma}

\begin{proof}
Multiply \eqref{eq:canonical-coordinate} by \(\overline{v_j}\) and take imaginary parts.  Since
\[
\Im(\overline{v_j}v_{j-1})=-\mathcal J_{j-1},
\qquad
\Im(\overline{v_j}v_{j+1})=\mathcal J_j,
\]
we obtain \eqref{eq:current-identity}.

For the derivative, the singular-value perturbation formula gives
\[
\partial_y s(x,y)
=
\Re(u^*\ii v)
=
-\Im(u^*v).
\]
Since \(u=J\overline v\),
\[
u^*v
=
v^{\mathsf T}Jv
=
\sum_{j=1}^n v_jv_{n+1-j}.
\]
Summing \eqref{eq:current-identity} over \(j=1,\dots,n\) gives
\[
(1-a)\sum_{j=1}^{n-1}\mathcal J_j
=
y+s(x,y)\Im(u^*v).
\]
Thus
\[
s(x,y)\partial_y s(x,y)
=
y-(1-a)\sum_{j=1}^{n-1}\mathcal J_j.
\]
Since
\[
\lambda_1(x,y)=s(x,y)^2,
\]
we get \eqref{eq:lambda-y-current}.
\end{proof}

\begin{lemma}[A real Wronskian invariant]\label{lem:real-wronskian}
Assume \(n\ge3\), \(y>0\), and let \(v\) be in canonical gauge.  Define
\[
W_j:=v_{j+1}v_{n+1-j}-av_jv_{n-j},
\qquad
j=0,\dots,n,
\]
with \(v_0=v_{n+1}=0\).  Then
\[
W_0=0,
\qquad
W_j-W_{j-1}
=
s(x,y)\bigl(|v_j|^2-|v_{n+1-j}|^2\bigr)\in\R.
\]
Consequently,
\[
W_j\in\R,
\qquad
j=0,\dots,n.
\]
\end{lemma}

\begin{proof}
Using \eqref{eq:canonical-coordinate} at the indices \(j\) and \(n+1-j\), we have
\[
v_{j+1}+av_{j-1}
=
(x+\ii y)v_j-s(x,y)\overline{v_{n+1-j}},
\]
and
\[
v_{n+2-j}+av_{n-j}
=
(x+\ii y)v_{n+1-j}-s(x,y)\overline{v_j}.
\]
Therefore
\[
W_j-W_{j-1}
=
v_{n+1-j}(v_{j+1}+av_{j-1})
-
v_j(v_{n+2-j}+av_{n-j})
=
s(x,y)\bigl(|v_j|^2-|v_{n+1-j}|^2\bigr).
\]
Since \(W_0=0\), all \(W_j\) are real.
\end{proof}

\begin{lemma}[Imaginary-axis seed]\label{lem:imaginary-axis-seed}
Let \(n\ge3\) and \(y>0\).  At \(x=0\), the canonical ground-state vector \(v\) may be written in the form
\[
v_j=\beta\,\ii^{\,j-1}w_j,
\qquad
w_j>0,
\qquad
j=1,\dots,n,
\]
for some unimodular \(\beta\).  Moreover,
\[
\Im\left(\frac{v_{j+1}}{v_j}\right)>0,
\qquad
j=1,\dots,n-1,
\]
and
\[
\Im\left(\frac{\overline{v_{n+1-j}}}{v_j}\right)>0,
\qquad
j=1,\dots,n.
\]
\end{lemma}

\begin{proof}
Let
\[
U:=\diag(1,\ii,\ii^2,\dots,\ii^{\,n-1}).
\]
A direct computation shows that
\[
U^*\mathcal H_0(y)U
\]
is a real symmetric irreducible \(Z\)-matrix: its first off-diagonal entries are
\[
-(1-a)y<0,
\]
and its second off-diagonal entries are
\[
-a<0.
\]
Choosing \(\Lambda\) larger than the largest diagonal entry, the matrix
\[
\Lambda I_n-U^*\mathcal H_0(y)U
\]
is entrywise nonnegative and irreducible.  By Perron--Frobenius, the lowest eigenvalue of \(U^*\mathcal H_0(y)U\) is simple and has a strictly positive eigenvector
\[
w=(w_1,\dots,w_n)^T,
\qquad
w_j>0.
\]
Thus a ground-state vector of \(\mathcal H_0(y)\) has the form
\[
v_j=\beta\,\ii^{\,j-1}w_j.
\]

It remains to use the canonical gauge.  Insert this form into
\[
\ii yv_j-av_{j-1}-v_{j+1}
=
s(0,y)\overline{v_{n+1-j}}.
\]
After division by \(\beta\ii^j\), one obtains
\[
w_{j+1}-yw_j-aw_{j-1}
=
-\gamma s(0,y)w_{n+1-j},
\qquad
\gamma:=\frac{\overline\beta}{\beta}(-\ii)^n.
\]
For \(j=n\), this gives
\[
yw_n+aw_{n-1}
=
\gamma s(0,y)w_1.
\]
The left-hand side is strictly positive, so
\[
\gamma=1.
\]
Consequently,
\[
\frac{v_{j+1}}{v_j}
=
\ii\,\frac{w_{j+1}}{w_j},
\]
and hence
\[
\Im\left(\frac{v_{j+1}}{v_j}\right)>0.
\]
Also,
\[
\frac{\overline{v_{n+1-j}}}{v_j}
=
\ii\,\frac{w_{n+1-j}}{w_j},
\]
so
\[
\Im\left(\frac{\overline{v_{n+1-j}}}{v_j}\right)>0.
\]
\end{proof}

\begin{lemma}[Angular identity]\label{lem:angular-identity}
Assume \(n\ge3\), \(y>0\), and let \(v\) be in canonical gauge.  Since all entries of \(v\) are nonzero by Lemmas \ref{lem:no-interior-zero} and \ref{lem:no-boundary-zero}, define
\[
r_j:=\frac{v_{j+1}}{v_j},
\qquad
j=1,\dots,n-1,
\]
and set
\[
r_0:=\infty,
\qquad
r_n:=0.
\]
For \(j=1,\dots,n-1\), put
\[
p_j:=\frac{\overline{v_{n+1-j}}}{v_j},
\qquad
\alpha_j:=(x+\ii y)-\frac{a}{r_{j-1}},
\qquad
\beta_j:=(x-\ii y)-\overline{r_{n+1-j}}.
\]
Then
\[
\Im(\alpha_j\overline{p_j})
=
-\Im(\beta_jp_j).
\]
\end{lemma}

\begin{proof}
From \eqref{eq:canonical-coordinate},
\[
\alpha_jv_j
=
v_{j+1}+s(x,y)\overline{v_{n+1-j}}.
\]
Multiplying by \(v_{n+1-j}/|v_j|^2\) gives
\[
\alpha_j\overline{p_j}
=
\frac{v_{j+1}v_{n+1-j}+s(x,y)|v_{n+1-j}|^2}{|v_j|^2},
\]
and hence
\[
\Im(\alpha_j\overline{p_j})
=
\frac{\Im(v_{j+1}v_{n+1-j})}{|v_j|^2}.
\]

Now conjugate \eqref{eq:canonical-coordinate} at the index \(n+1-j\).  This gives
\[
(x-\ii y)\overline{v_{n+1-j}}
-a\overline{v_{n-j}}
-\overline{v_{n+2-j}}
=
s(x,y)v_j.
\]
Dividing by \(v_j\) yields
\[
\beta_jp_j
=
s(x,y)+a\frac{\overline{v_{n-j}}}{v_j}.
\]
Thus
\[
\Im(\beta_jp_j)
=
-\frac{a\Im(v_jv_{n-j})}{|v_j|^2}.
\]
By Lemma \ref{lem:real-wronskian},
\[
\Im(v_{j+1}v_{n+1-j})
=
a\Im(v_jv_{n-j}),
\]
and the claim follows.
\end{proof}

\begin{lemma}[Phase continuation]\label{lem:phase-continuation}
Fix \(y>0\).  On any open interval in \(x\) on which the canonical vector \(v(x)\) is real analytic and has no zero components, if at one point \(x_0\) one has
\[
\Im\left(\frac{v_{j+1}(x_0)}{v_j(x_0)}\right)>0,
\qquad
j=1,\dots,n-1,
\]
and
\[
\Im\left(\frac{\overline{v_{n+1-j}(x_0)}}{v_j(x_0)}\right)>0,
\qquad
j=1,\dots,n,
\]
then the same inequalities hold throughout the interval.  Consequently,
\[
\partial_y\lambda_1(x,y)\ge0
\]
throughout the interval.
\end{lemma}

\begin{proof}
Write
\[
v_j(x)=t_j(x)e^{\ii\phi_j(x)},
\qquad
t_j(x)>0,
\]
with continuous arguments \(\phi_j\), and set
\[
\delta_j(x):=\phi_{j+1}(x)-\phi_j(x),
\qquad
j=1,\dots,n-1.
\]
Then
\[
\Im\left(\frac{v_{j+1}}{v_j}\right)>0
\]
is equivalent to
\[
\delta_j\in(0,\pi)
\]
modulo \(2\pi\), after choosing continuous representatives.

Write
\[
c=(1+a)x-\ii(1-a)y=|c|e^{-\ii\vartheta},
\qquad
0<\vartheta<\pi.
\]
For fixed amplitudes \(t_1,\dots,t_n\), the phase-dependent part of the Rayleigh quotient \(\langle v,\mathcal H_x(y)v\rangle\) is
\[
\mathcal P
=
-2|c|\sum_{j=1}^{n-1}t_jt_{j+1}\cos(\delta_j-\vartheta)
+
2a\sum_{j=1}^{n-2}t_jt_{j+2}\cos(\delta_j+\delta_{j+1}).
\]
Because \(v\) is a ground-state vector, its phase vector minimizes \(\mathcal P\) for the fixed amplitudes.

Let \(E_0\) be the subset of the interval on which
\[
\delta_j\in(0,\pi)
\qquad
(j=1,\dots,n-1).
\]
By assumption, \(E_0\) is nonempty.  It is open by continuity.  We prove it is closed.  Let \(x_*\in\overline E_0\).  Then
\[
\delta_j(x_*)\in[0,\pi].
\]
Suppose \(\delta_k(x_*)=0\) or \(\delta_k(x_*)=\pi\).  Vary the tail phases by
\[
\phi_m\mapsto\phi_m+\varepsilon,
\qquad
m\ge k+1.
\]
This changes only \(\delta_k\), replacing it by \(\delta_k+\varepsilon\).  The derivative of \(\mathcal P\) at \(\varepsilon=0\) is
\[
2|c|t_kt_{k+1}\sin(\delta_k-\vartheta)
-2a\mathbf 1_{\{k\ge2\}}t_{k-1}t_{k+1}\sin(\delta_{k-1}+\delta_k)
-2a\mathbf 1_{\{k\le n-2\}}t_kt_{k+2}\sin(\delta_k+\delta_{k+1}).
\]
If \(\delta_k=0\), this derivative is strictly negative.  If \(\delta_k=\pi\), it is strictly positive.  Either case contradicts minimality of the phase vector.  Hence no \(\delta_k\) can hit \(0\) or \(\pi\), so \(E_0\) is closed.  Since the interval is connected, \(E_0\) is the whole interval.  Therefore
\[
\Im\left(\frac{v_{j+1}}{v_j}\right)>0,
\qquad
j=1,\dots,n-1.
\]

Next set
\[
p_j:=\frac{\overline{v_{n+1-j}}}{v_j}.
\]
We show that
\[
\Im p_j>0,
\qquad
j=1,\dots,n.
\]
Suppose, for some \(1\le j\le n-1\), that \(p_j\) is real at a point \(x_*\).  Since all entries of \(v\) are nonzero, \(p_j(x_*)\ne0\).  By the first part of the proof,
\[
\Im\left((x+\ii y)-\frac{a}{r_{j-1}}\right)>0.
\]
Using Lemma \ref{lem:angular-identity}, and using that \(p_j(x_*)\) is real, we get
\[
\Im\left((x-\ii y)-\overline{r_{n+1-j}}\right)<0.
\]
The conjugated singular-vector equation gives
\[
\overline{r_{n-j}}
=
\frac{a}
{(x-\ii y)-\overline{r_{n+1-j}}-s(x,y)/p_j}.
\]
The denominator has negative imaginary part, because \(s(x,y)/p_j\in\R\).  Therefore
\[
\Im\overline{r_{n-j}}>0,
\]
or equivalently
\[
\Im r_{n-j}<0,
\]
contradicting the already proved inequality
\[
\Im r_{n-j}>0.
\]
Thus
\[
\Im p_j>0,
\qquad
j=1,\dots,n-1.
\]
Finally,
\[
p_n=\frac{1}{\overline{p_1}},
\]
so \(\Im p_n>0\) as well.

Since
\[
p_j
=
\frac{\overline{v_{n+1-j}}}{v_j}
=
\frac{\overline{v_jv_{n+1-j}}}{|v_j|^2},
\]
the inequality \(\Im p_j>0\) implies
\[
\Im(v_jv_{n+1-j})<0.
\]
The current identity \eqref{eq:current-identity} therefore gives
\[
\mathcal J_j-a\mathcal J_{j-1}
\le
y|v_j|^2,
\qquad
j=1,\dots,n.
\]
By induction,
\[
\mathcal J_j
\le
y\sum_{k=1}^j a^{j-k}|v_k|^2,
\qquad
j=1,\dots,n-1.
\]
Summing over \(j\) gives
\[
(1-a)\sum_{j=1}^{n-1}\mathcal J_j
\le
y\sum_{k=1}^{n-1}(1-a^{n-k})|v_k|^2
\le
y.
\]
Using \eqref{eq:lambda-y-current}, we obtain
\[
\partial_y\lambda_1(x,y)
=
2y-2(1-a)\sum_{j=1}^{n-1}\mathcal J_j
\ge0.
\]
\end{proof}

\begin{proof}[Proof of Theorem \ref{thm:vertical}]
For \(n=2\), the result follows from Lemma \ref{lem:vertical-n2}.  Assume \(n\ge3\).

Fix \(y_*>0\).  By Corollary \ref{cor:lowest-simple}, the lowest eigenvalue
\[
\lambda_1(x,y_*)
\]
is simple for every \(x\in\R\).  Hence Lemma \ref{lem:canonical-gauge} gives a real-analytic canonical ground-state vector \(v(x)\) on \(\R\).  Lemmas \ref{lem:no-interior-zero} and \ref{lem:no-boundary-zero} show that no component of \(v(x)\) vanishes.

At \(x=0\), Lemma \ref{lem:imaginary-axis-seed} gives
\[
\Im\left(\frac{v_{j+1}(0)}{v_j(0)}\right)>0,
\qquad
j=1,\dots,n-1,
\]
and
\[
\Im\left(\frac{\overline{v_{n+1-j}(0)}}{v_j(0)}\right)>0,
\qquad
j=1,\dots,n.
\]
Applying Lemma \ref{lem:phase-continuation} on the interval \(\R\), we obtain
\[
\partial_y\lambda_1(x,y_*)\ge0,
\qquad
x\in\R.
\]
Since \(y_* >0\) was arbitrary,
\[
\partial_y\lambda_1(x,y)\ge0,
\qquad
x\in\R,\ y>0.
\]

Because \(A_n(a)\) is real,
\[
s_{\min}((x-\ii y)I_n-A_n(a))
=
s_{\min}((x+\ii y)I_n-A_n(a)),
\]
so \(\lambda_1(x,y)\) is even in \(y\).  Hence \(y\mapsto\lambda_1(x,y)\) is nondecreasing on \([0,\infty)\).  Since
\[
s_{\min}((x+\ii y)I_n-A_n(a))=\sqrt{\lambda_1(x,y)},
\]
the same is true for the least singular value.  This proves
\[
s_{\min}((x+\ii y)I_n-A_n(a))
\ge
s_{\min}(xI_n-A_n(a)).
\]
\end{proof}

\begin{proposition}[Reduction to the real trace]\label{prop:planar-reduction}
For every \(0<a<1\), \(n\ge2\), and \(\eps>0\),
\[
\sigma_\eps(A_n(a))\ \text{is connected}
\quad\Longleftrightarrow\quad
E_{n,\eps}(a)\ \text{is connected}.
\]
Consequently,
\[
\sigma_\eps(A_n(a))\ \text{is connected}
\quad\Longleftrightarrow\quad
\eps>\tau_n(a).
\]
Moreover, every connected component of \(\sigma_\eps(A_n(a))\) is contractible.  In particular, if \(\sigma_\eps(A_n(a))\) is connected, then it is simply connected.
\end{proposition}

\begin{proof}
Let
\[
\Sigma:=\sigma_\eps(A_n(a)).
\]
For fixed \(x\in\R\), define the vertical section
\[
\Omega_x:=\{y\in\R:x+\ii y\in\Sigma\}.
\]
By Theorem \ref{thm:vertical} and the symmetry \(s(x,-y)=s(x,y)\), the set \(\Omega_x\) is either empty or an open interval symmetric about \(0\).  It is nonempty if and only if
\[
s_{\min}(xI_n-A_n(a))<\eps,
\]
that is, if and only if \(x\in E_{n,\eps}(a)\).  Hence the vertical projection
\[
\pi:\Sigma\to\R,
\qquad
\pi(x+\ii y)=x,
\]
satisfies
\[
\pi(\Sigma)=E_{n,\eps}(a).
\]

If \(\Sigma\) is connected, then \(E_{n,\eps}(a)\) is connected as the continuous image of a connected set.

Conversely, suppose \(E_{n,\eps}(a)\) is connected.  Let
\[
z_0=x_0+\ii y_0,
\qquad
z_1=x_1+\ii y_1
\]
be two points of \(\Sigma\).  By vertical monotonicity, the vertical segments from \(z_0\) to \(x_0\), and from \(x_1\) to \(z_1\), lie in \(\Sigma\).  Since \(E_{n,\eps}(a)\) is an interval, the horizontal segment from \(x_0\) to \(x_1\) lies in
\[
E_{n,\eps}(a)\subset\Sigma.
\]
Concatenating these three segments gives a path in \(\Sigma\) from \(z_0\) to \(z_1\).  Thus \(\Sigma\) is path-connected, hence connected.

The equivalence with
\[
\eps>\tau_n(a)
\]
now follows from Proposition \ref{prop:real-axis-connectedness}.

It remains to prove the contractibility of components.  Let \(\Omega\) be a connected component of \(\Sigma\).  If \(x+\ii y\in\Omega\), then the vertical segment from \(x+\ii y\) to \(x\) lies in \(\Sigma\), hence lies in the same connected component \(\Omega\).  Thus every component meets the real axis.  Moreover, if \(x\in\Omega\cap\R\), then the whole vertical section
\[
\{x+\ii y:x+\ii y\in\Sigma\}
\]
lies in \(\Omega\), again by vertical monotonicity.  Therefore
\[
\Omega
=
\{x+\ii y:\ x\in I_\Omega,\ |y|<h_\Omega(x)\},
\]
where
\[
I_\Omega:=\Omega\cap\R
\]
is a connected component of \(E_{n,\eps}(a)\), hence an open interval, and \(h_\Omega(x)>0\) is the vertical half-height of the section.

The homotopy
\[
\Phi_t(x+\ii y):=x+\ii(1-t)y,
\qquad
0\le t\le1,
\]
stays inside \(\Omega\).  Hence \(\Omega\) deformation retracts onto the interval \(I_\Omega\).  Therefore \(\Omega\) is contractible.
\end{proof}

\section{The finite monotonicity conjecture and threshold consequences}

\begin{conjecture}[Real-axis barrier monotonicity]\label{conj:barrier}
For every \(0<a<1\) and every \(n\ge2\),
\[
\tau_{n+1}(a)\le\tau_n(a).
\]
\end{conjecture}

\begin{remark}
Conjecture \ref{conj:barrier} is the remaining finite comparison principle needed to turn the exact connectedness criterion into a first-hit equals eventual-hit theorem.  It is not a consequence of ordinary Hermitian eigenvalue interlacing, because \(A_n(a)\) is only nonunitarily similar to a Hermitian matrix and the Euclidean singular values are not preserved by this similarity.
\end{remark}

For fixed \(a\in(0,1)\) and \(\eps>0\), define
\[
N_f(a,\eps)
:=
\min\{n\ge2:\sigma_\eps(A_n(a))\ \text{is connected}\},
\]
and
\[
N_e(a,\eps)
:=
\min\{N\ge2:\sigma_\eps(A_n(a))\ \text{is connected for every }n\ge N\}.
\]

\begin{theorem}[Threshold identity under barrier monotonicity]\label{thm:conditional-threshold}
Assume Conjecture \ref{conj:barrier}.  Then
\[
N_f(a,\eps)=N_e(a,\eps).
\]
Writing the common value as
\[
N_c(a,\eps):=N_f(a,\eps)=N_e(a,\eps),
\]
one has
\[
N_c(a,\eps)
=
\min\{n\ge2:\eps>\tau_n(a)\}.
\]
\end{theorem}

\begin{proof}
By Proposition \ref{prop:planar-reduction},
\[
\sigma_\eps(A_n(a))\ \text{is connected}
\quad\Longleftrightarrow\quad
\eps>\tau_n(a).
\]
By Conjecture \ref{conj:barrier}, the sequence \(\tau_n(a)\) is nonincreasing.  The upper bound in Theorem \ref{thm:tau-bounds} below implies
\[
\tau_n(a)\le(n+1)a^{n/2}\to0.
\]
Hence the set
\[
\{n\ge2:\eps>\tau_n(a)\}
\]
is a nonempty tail set.  Its first element is both the first connectedness index and the eventual connectedness index.
\end{proof}

\section{Explicit Green kernel}

The quantitative estimates use the inverse of
\[
B_n(x)=xI_n-A_n(a)
\]
for \(x\) away from the spectrum.  Write
\[
q:=\sqrt a.
\]
For
\[
x=2q\cos\theta,
\qquad
\sin((n+1)\theta)\ne0,
\]
the matrix \(B_n(x)\) is invertible.

\begin{lemma}[Green kernel]\label{lem:green}
Let \(N:=n+1\), \(q:=\sqrt a\), and \(x=2q\cos\theta\), with \(\sin\theta\ne0\) and \(\sin(N\theta)\ne0\).  Then
\[
(B_n(x)^{-1})_{ij}
=
\frac{
q^{i-j-1}\sin(i\theta)\sin((N-j)\theta)
}{
\sin\theta\,\sin(N\theta)
},
\qquad i\le j,
\]
and
\[
(B_n(x)^{-1})_{ij}
=
\frac{
q^{i-j-1}\sin(j\theta)\sin((N-i)\theta)
}{
\sin\theta\,\sin(N\theta)
},
\qquad i>j.
\]
Only absolute values will be used; the harmless global sign depends on the convention for the inverse of the Dirichlet Jacobi matrix.
\end{lemma}

\begin{proof}
Let
\[
\Delta_n=\diag(1,q,q^2,\dots,q^{n-1}).
\]
Then
\[
\Delta_n^{-1}A_n(a)\Delta_n=qT_n,
\]
where
\[
T_n=\tridiag(1,0,1)
\]
is the symmetric Dirichlet Jacobi matrix.  Hence
\[
B_n(x)
=
\Delta_n(xI_n-qT_n)\Delta_n^{-1}.
\]
Thus
\[
B_n(x)^{-1}
=
\Delta_n(xI_n-qT_n)^{-1}\Delta_n^{-1}.
\]
For \(x=2q\cos\theta\),
\[
xI_n-qT_n
=
q(2\cos\theta\,I_n-T_n).
\]
The standard Dirichlet Green kernel for \(2\cos\theta\,I_n-T_n\) is
\[
\left(2\cos\theta\,I_n-T_n\right)^{-1}_{ij}
=
\frac{
\sin(\min\{i,j\}\theta)\sin((N-\max\{i,j\})\theta)
}{
\sin\theta\,\sin(N\theta)
},
\]
up to an overall sign.  Multiplying by \(q^{-1}\) and by the diagonal similarity factor \(q^{i-j}\) gives the displayed formulas.
\end{proof}

\section{Sharp two-sided bounds with \texorpdfstring{\(S_2\)}{S2}}

Define
\[
S_2(a):=\sum_{m=1}^{\infty}m^2a^m.
\]
The closed form is
\[
S_2(a)=\frac{a(1+a)}{(1-a)^3}.
\]
Set
\[
c_2(a)
:=
\frac{\sqrt8\,a}{3\pi^2S_2(a)}
=
\frac{\sqrt8(1-a)^3}{3\pi^2(1+a)}.
\]

\begin{theorem}[Two-sided real-axis barrier bounds]\label{thm:tau-bounds}
For every \(0<a<1\) and every \(n\ge2\),
\[
c_2(a)(n+1)a^{n/2}
\le
\tau_n(a)
\le
(n+1)a^{n/2}.
\]
\end{theorem}

\begin{proof}
We first prove the upper bound.  Let
\[
x=2\sqrt a\cos\theta,
\qquad 0<\theta<\pi.
\]
Define
\[
q^{(n)}(\theta)
=
\left(
\sin\theta,\,
a^{1/2}\sin2\theta,\,
a\sin3\theta,\,
\dots,\,
a^{(n-1)/2}\sin n\theta
\right)^T.
\]
Using
\[
2\cos\theta\sin(j\theta)
=
\sin((j+1)\theta)+\sin((j-1)\theta),
\]
one obtains
\[
B_n(x)q^{(n)}(\theta)
=
a^{n/2}\sin((n+1)\theta)e_n.
\]
Therefore
\[
\mu_n(x;a)
\le
\frac{
a^{n/2}|\sin((n+1)\theta)|
}{
\left(\sum_{j=1}^n a^{j-1}\sin^2(j\theta)\right)^{1/2}
}.
\]
Since
\[
\left(\sum_{j=1}^n a^{j-1}\sin^2(j\theta)\right)^{1/2}
\ge |\sin\theta|
\]
and
\[
|\sin((n+1)\theta)|\le(n+1)|\sin\theta|,
\]
we get
\[
\mu_n(x;a)\le(n+1)a^{n/2}
\]
for every \(x\in(-2\sqrt a,2\sqrt a)\).  All real spectral gaps lie in this interval, hence
\[
\tau_n(a)\le(n+1)a^{n/2}.
\]

We now prove the lower bound.  Let
\[
N:=n+1,
\qquad
\delta:=\frac{3\pi}{2N},
\qquad
\theta:=\pi-\delta.
\]
Then
\[
\frac{(n-1)\pi}{n+1}<\theta<\frac{n\pi}{n+1}.
\]
Consequently,
\[
x_*:=2\sqrt a\cos\theta
\]
lies in the leftmost real spectral gap.  Thus
\[
\tau_n(a)\ge \mu_n(x_*;a).
\]
Since \(x_*\notin\sigma(A_n(a))\),
\[
\mu_n(x_*;a)
=
\frac1{\|B_n(x_*)^{-1}\|_2}
\ge
\frac1{\|B_n(x_*)^{-1}\|_F}.
\]

We estimate the Frobenius norm using Lemma \ref{lem:green}.  For \(i\le j\), set
\[
p:=i,
\qquad
q_0:=N-j,
\qquad
m:=p+q_0.
\]
Then \(p,q_0\ge1\), \(m\le N\), and
\[
i-j-1=m-N-1.
\]
Moreover,
\[
|\sin(N\theta)|=1,
\qquad
\sin\theta=\sin\delta,
\]
and, since \(0<\delta\le\pi/2\),
\[
\sin\delta\ge \frac{2\delta}{\pi}.
\]
Also,
\[
|\sin(r\theta)|=|\sin(r\delta)|\le r\delta.
\]
Hence for \(i\le j\),
\[
|(B_n(x_*)^{-1})_{ij}|^2
\le
\frac{\pi^2}{4}\,
p^2q_0^2\delta^2\,
a^{m-N-1}.
\]
For \(i>j\), the same upper bound holds after enlarging the estimate, because the corresponding exponent is \(N-m-1\), and
\[
a^{N-m-1}\le a^{m-N-1}
\qquad(0<a<1,\ m\le N).
\]
Therefore
\[
\|B_n(x_*)^{-1}\|_F^2
\le
\frac{\pi^2}{2}\delta^2a^{-N-1}
\sum_{\substack{p,q_0\ge1\\p+q_0\le N}}
p^2q_0^2a^{p+q_0}.
\]
Extending the sum to all \(p,q_0\ge1\) gives
\[
\|B_n(x_*)^{-1}\|_F^2
\le
\frac{\pi^2}{2}\delta^2a^{-N-1}
\left(\sum_{p=1}^{\infty}p^2a^p\right)^2.
\]
Thus
\[
\|B_n(x_*)^{-1}\|_F^2
\le
\frac{\pi^2}{2}\delta^2a^{-N-1}S_2(a)^2.
\]
Since
\[
\delta=\frac{3\pi}{2N},
\]
we obtain
\[
\|B_n(x_*)^{-1}\|_F
\le
\frac{3\pi^2}{\sqrt8}\,
S_2(a)\,
\frac{a^{-(N+1)/2}}{N}.
\]
Therefore
\[
\mu_n(x_*;a)
\ge
\frac{\sqrt8}{3\pi^2S_2(a)}Na^{(N+1)/2}.
\]
Since \(N=n+1\),
\[
Na^{(N+1)/2}
=
(n+1)a^{(n+2)/2}
=
a(n+1)a^{n/2}.
\]
Hence
\[
\tau_n(a)
\ge
\frac{\sqrt8\,a}{3\pi^2S_2(a)}
(n+1)a^{n/2}
=
c_2(a)(n+1)a^{n/2}.
\]
The proof is complete.
\end{proof}

\section{Explicit threshold bounds}

Define the first-hit connectedness threshold by
\[
N_c(a,\eps):=\min\{n\ge2:\sigma_\eps(A_n(a))\ \text{is connected}\}.
\]
By Proposition \ref{prop:planar-reduction},
\[
N_c(a,\eps)=\min\{n\ge2:\eps>\tau_n(a)\}.
\]
This equality is a first-hit formula.  If Conjecture \ref{conj:barrier} holds, it is also the eventual threshold.

Define
\[
N_+(a,\eps)
:=
\min\{n\ge2:(n+1)a^{n/2}<\eps\},
\]
and
\[
N_-(a,\eps)
:=
1+\max\{n\ge2:c_2(a)(n+1)a^{n/2}\ge\eps\},
\]
with the convention that the maximum over the empty set is \(1\).

\begin{corollary}[Two-sided finite-size threshold bounds]\label{cor:Nc-bounds}
For every \(0<a<1\) and every \(\eps>0\),
\[
N_-(a,\eps)\le N_c(a,\eps)\le N_+(a,\eps).
\]
If Conjecture \ref{conj:barrier} holds, the same bounds apply to the common threshold
\[
N_e(a,\eps)=N_f(a,\eps).
\]
\end{corollary}

\begin{proof}
If
\[
(n+1)a^{n/2}<\eps,
\]
then Theorem \ref{thm:tau-bounds} gives
\[
\tau_n(a)<\eps.
\]
Hence \(\sigma_\eps(A_n(a))\) is connected, and so
\[
N_c(a,\eps)\le n.
\]
Taking the least such \(n\) gives
\[
N_c(a,\eps)\le N_+(a,\eps).
\]

Conversely, if
\[
c_2(a)(n+1)a^{n/2}\ge\eps,
\]
then Theorem \ref{thm:tau-bounds} gives
\[
\tau_n(a)\ge\eps.
\]
Thus \(\sigma_\eps(A_n(a))\) is not connected.  Therefore the connectedness threshold must be larger than every such \(n\), which gives
\[
N_c(a,\eps)\ge N_-(a,\eps).
\]
\end{proof}

Let
\[
\gamma:=\frac{|\log a|}{2}.
\]
Then
\[
a^{n/2}=e^{-\gamma n}.
\]
The equation
\[
C(n+1)e^{-\gamma n}=\eps
\]
with \(C>0\) is solved by the negative branch \(W_{-1}\) of Lambert's \(W\)-function:
\[
n
=
-\frac1\gamma
W_{-1}\!\left(-\frac{\gamma\eps e^{-\gamma}}{C}\right)-1.
\]

\begin{corollary}[Small-\(\eps\) scale]\label{cor:small-eps}
For fixed \(0<a<1\),
\[
N_c(a,\eps)
=
\frac{2}{|\log a|}
\left(
\log\frac1\eps+\log\log\frac1\eps
\right)
+
O_a(1),
\qquad
\eps\downarrow0,
\]
for the first-hit threshold.  If Conjecture \ref{conj:barrier} holds, the same asymptotic applies to the common eventual/first-hit threshold.
\end{corollary}

\begin{proof}
Apply Corollary \ref{cor:Nc-bounds} with \(C=1\) and \(C=c_2(a)\).  The standard expansion
\[
W_{-1}(-t)
=
\log t-\log|\log t|+O\!\left(\frac{\log|\log t|}{|\log t|}\right),
\qquad t\downarrow0,
\]
gives
\[
N_c(a,\eps)
=
\frac1\gamma
\left(
\log\frac1\eps+\log\log\frac1\eps
\right)
+
O_a(1).
\]
Since \(\gamma=|\log a|/2\), the result follows.
\end{proof}

\section{Physical interpretation: the open-boundary Hatano--Nelson chain}

The matrix \(A_n(a)\) is the finite open-boundary Hatano--Nelson Hamiltonian with asymmetric nearest-neighbor hopping after the normalization
\[
t_R=1,
\qquad
t_L=a,
\qquad 0<a<1.
\]
Writing
\[
a=e^{-2g},
\qquad g>0,
\]
one has
\[
\Delta_n^{-1}A_n(a)\Delta_n
=
e^{-g}\tridiag(1,0,1),
\qquad
\Delta_n=\diag(1,e^{-g},e^{-2g},\dots,e^{-(n-1)g}).
\]
This is the open-boundary imaginary-gauge transformation.  It removes the hopping nonreciprocity at the spectral level but is nonunitary, and therefore does not preserve Euclidean pseudospectra.

The right eigenvectors have the form
\[
r^{(j)}_\ell
=
a^{(\ell-1)/2}\sin\frac{j\ell\pi}{n+1},
\qquad
\ell=1,\dots,n,
\]
whereas the left eigenvectors have the reciprocal profile
\[
\ell^{(j)}_\ell
=
a^{-(\ell-1)/2}\sin\frac{j\ell\pi}{n+1}.
\]
Thus right and left eigenvectors are exponentially biased toward opposite edges.  This is the finite-chain non-Hermitian skin effect.

The connectedness threshold has a direct physical interpretation.  At finite resolution \(\eps\), the open-chain pseudospectrum initially consists of separated spectral islands around individual open-boundary eigenvalues.  The threshold is the length scale at which the largest real-axis barrier is below the resolution level, so the islands merge.  The two-sided estimate
\[
c_2(a)(n+1)a^{n/2}
\le
\tau_n(a)
\le
(n+1)a^{n/2}
\]
shows that the barrier is exponentially small in the chain length, with exponent controlled by the nonreciprocity parameter \(g=|\log a|/2\).  Accordingly,
\[
N_c(e^{-2g},\eps)
=
\frac1g
\left(
\log\frac1\eps+\log\log\frac1\eps
\right)
+
O_g(1).
\]
Stronger nonreciprocity lowers the chain length needed for pseudospectral connectedness.  In the reciprocal limit \(a\uparrow1\), the scale diverges, consistent with the Hermitian limit.

\section{Concluding remarks}

The rigorous unconditional contribution of this paper is the vertical monotonicity theorem, the exact real-axis barrier formulation of planar connectedness, and the sharp two-sided estimate
\[
\tau_n(a)\asymp_a (n+1)a^{n/2}.
\]
The use of \(S_2(a)\) follows from summing the Green-kernel factors exactly:
\[
\sum_{p,q\ge1}p^2q^2a^{p+q}
=
\left(\sum_{p=1}^{\infty}p^2a^p\right)^2
=
S_2(a)^2.
\]
The finite monotonicity conjecture
\[
\tau_{n+1}(a)\le\tau_n(a)
\]
is the remaining comparison principle needed for the unconditional identity
\[
N_e(a,\eps)=N_f(a,\eps).
\]
A proof of this inequality would complete the finite-dimensional first-hit equals eventual-hit theorem for this nonnormal Toeplitz family.

\begin{thebibliography}{99}

\bibitem{AmmariBarandunCaoDaviesHiltunen2024}
H.~Ammari, S.~Barandun, J.~Cao, B.~Davies, and E.~O. Hiltunen,
Mathematical foundations of the non-Hermitian skin effect,
\emph{Arch. Ration. Mech. Anal.} \textbf{248} (2024), article no.~33.

\bibitem{BottcherGrudsky2005}
A.~B\"ottcher and S.~M. Grudsky,
\emph{Spectral Properties of Banded Toeplitz Matrices},
SIAM, Philadelphia, 2005.

\bibitem{BottcherEmbreeTrefethen2002}
A.~B\"ottcher, M.~Embree, and L.~N. Trefethen,
Piecewise continuous Toeplitz matrices and operators: slow approach to infinity,
\emph{SIAM J. Matrix Anal. Appl.} \textbf{24} (2002), no.~2, 484--489.

\bibitem{CorlessGonnetHareJeffreyKnuth1996}
R.~M. Corless, G.~H. Gonnet, D.~E.~G. Hare, D.~J. Jeffrey, and D.~E. Knuth,
On the Lambert \(W\) function,
\emph{Adv. Comput. Math.} \textbf{5} (1996), 329--359.

\bibitem{HatanoNelson1996}
N.~Hatano and D.~R. Nelson,
Localization transitions in non-Hermitian quantum mechanics,
\emph{Phys. Rev. Lett.} \textbf{77} (1996), no.~3, 570--573.

\bibitem{HornJohnson2013}
R.~A. Horn and C.~Johnson,
\emph{Matrix Analysis},
2nd ed.,
Cambridge University Press, Cambridge, 2013.

\bibitem{KunstEdvardssonBudichBergholtz2018}
F.~K. Kunst, E.~Edvardsson, J.~C. Budich, and E.~J. Bergholtz,
Biorthogonal bulk-boundary correspondence in non-Hermitian systems,
\emph{Phys. Rev. Lett.} \textbf{121} (2018), 026808.

\bibitem{NoschesePasquiniReichel2013}
S.~Noschese, L.~Pasquini, and L.~Reichel,
Tridiagonal Toeplitz matrices: properties and novel applications,
\emph{Numer. Linear Algebra Appl.} \textbf{20} (2013), no.~2, 302--326.

\bibitem{OkumaKawabataShiozakiSato2020}
N.~Okuma, K.~Kawabata, K.~Shiozaki, and M.~Sato,
Topological origin of non-Hermitian skin effects,
\emph{Phys. Rev. Lett.} \textbf{124} (2020), 086801.

\bibitem{ReichelTrefethen1992}
L.~Reichel and L.~N. Trefethen,
Eigenvalues and pseudo-eigenvalues of Toeplitz matrices,
\emph{Linear Algebra Appl.} \textbf{162--164} (1992), 153--185.

\bibitem{Trefethen1997}
L.~N. Trefethen,
Pseudospectra of linear operators,
\emph{SIAM Rev.} \textbf{39} (1997), no.~3, 383--406.

\bibitem{TrefethenEmbree2005}
L.~N. Trefethen and M.~Embree,
\emph{Spectra and Pseudospectra: The Behavior of Nonnormal Matrices and Operators},
Princeton University Press, Princeton, 2005.

\bibitem{YaoWang2018}
S.~Yao and Z.~Wang,
Edge states and topological invariants of non-Hermitian systems,
\emph{Phys. Rev. Lett.} \textbf{121} (2018), 086803.

\end{thebibliography}

\end{document}